\medskip

\begin{enumerate}
\item  On a : $z_1=1-\i\sqrt{3}$.

Son module vaut : $|z_1|=\sqrt{1^2+\left(-\sqrt{3}\right)^2}=\sqrt{1+3}=2$.

On cherche un argument $\theta$ de $z_1$.

On a $\cos(\theta)=\dfrac{1}{2}$ et $\sin(\theta)=\dfrac{-\sqrt{3}}{2}$, on peut donc prendre $\theta=-\dfrac{\pi}{3}$.

Ainsi, la forme exponentielle de $z_1$ est : $z_1=2\e^{-\i\frac{\pi}{3}}$.

\item $z_1^2=\left(2\e^{-\i\frac{\pi}{3}}\right)^2=4\e^{-\i\frac{2\pi}{3}}$ et $z_2=4\e^{\i\frac{\pi}{6}}$.

Ainsi :
\begin{align*}
z &= z_1^2\times z_2 \\
&= 4\e^{-\i\frac{2\pi}{3}}\times 4\e^{\i\frac{\pi}{6}} \\
&= 16\e^{\i\left(-\frac{2\pi}{3}+\frac{\pi}{6}\right)} \\
&= 16\e^{-\i\frac{\pi}{2}} \\
&= 16\left(\cos\left(-\frac{\pi}{2}\right)+\i\sin\left(-\frac{\pi}{2}\right)\right) \\
&= 16(0+\i\times(-1)) \\
&= -16\i.
\end{align*}

Le nombre $z$ est donc bien un nombre imaginaire pur.
\end{enumerate}
